\chapter{Superfícies}\label{ch:b2:surfaces}

Depois das curvas, as superfícies: objetos a dois parâmetros em $\R^3$. O
cálculo diferencial de \cref{ch:b2:diffcalc} fornece tudo o
que precisamos --- as derivadas parciais dão vetores tangentes, o produto
vetorial dá a normal, os \omterm{def:b2:linalg:det}{determinantes} dão \omterm{def:b2:surfaces:area}{áreas}. Definimos as
\omterm{def:b2:surfaces:param}{superfícies parametrizadas} regulares, seus \omterm{def:b2:surfaces:tangent}{planos tangentes} e a
\emph{\omterm{def:b2:surfaces:fff}{primeira forma fundamental}}, que codifica todas as medidas de
\omterm{def:b2:curves:length}{comprimento} e de \omterm{def:b2:surfaces:area}{área} sobre a superfície. Superfícies também surgem como
conjuntos de nível $f(x, y, z) = c$; o gradiente então dirige a normal.

\section{Superfícies parametrizadas}

\begin{definition}[Superfície parametrizada regular]\label{def:b2:surfaces:param}
Seja $U \subseteq \R^2$ um \omterm{def:b2:metric:topology}{aberto}. Uma \emph{superfície
parametrizada}\index{superfície parametrizada} de
classe $\mathcal{C}^k$ ($k \geq 1$) é uma aplicação $\sigma \colon U \to
\R^3$, $(u, v) \mapsto \sigma(u, v)$, de classe $\mathcal{C}^k$. Um
ponto é \emph{regular} se os vetores derivadas parciais
\[
\sigma_u(u, v) = \frac{\partial\sigma}{\partial u}(u,v),
\qquad
\sigma_v(u, v) = \frac{\partial\sigma}{\partial v}(u,v)
\]
forem linearmente independentes, isto é, $\sigma_u \wedge \sigma_v \neq 0$;
a superfície é \emph{regular} se todos os seus pontos o forem.
\end{definition}

\begin{example}[As três descrições padrão]\label{ex:b2:surfaces:standard}
\leavevmode
\begin{enumerate}
\item \emph{Gráfico:} $\sigma(u, v) = (u,\ v,\ f(u, v))$ para $f \in
\mathcal{C}^1(U)$. Sempre regular: $\sigma_u = (1, 0, f_u)$ e
$\sigma_v = (0, 1, f_v)$ são independentes.
\item \emph{Esfera (coordenadas esféricas):} para a esfera de
raio $R$,
\[
\sigma(\theta, \varphi)
= (R\cos\theta\cos\varphi,\ R\sin\theta\cos\varphi,\
R\sin\varphi),
\qquad
(\theta, \varphi) \in \R \times
\bigl(-\tfrac\pi2, \tfrac\pi2\bigr),
\]
com $\theta$ a longitude e $\varphi$ a latitude. Verifica-se que
$\norm{\sigma_\theta \wedge \sigma_\varphi} = R^2\cos\varphi > 0$:
regular fora dos polos (que esta carta omite).
\item \emph{Conjunto de nível:} $S = \{(x,y,z) : f(x, y, z) = c\}$ em que
$f$ é $\mathcal{C}^1$ e $\nabla f \neq 0$ em $S$. Perto de cada
ponto, uma coordenada pode ser expressa como função das outras
duas pelo teorema da função implícita (\cref{ch:b2:diffcalc}), de modo que
$S$ é localmente um gráfico.
\end{enumerate}
\end{example}

\begin{example}[Do conjunto de nível ao gráfico]
\label{ex:b2:surfaces:levelgraph}
O teorema da função implícita do item 3 merece uma passagem
explícita. Tome a esfera $x^2 + y^2 + z^2 = R^2$ perto de seu polo
norte $(0, 0, R)$: ali $\frac{\partial f}{\partial z} = 2z =
2R \neq 0$ e, resolvendo em $z$, obtém-se a carta gráfico
\[
z = \sqrt{R^2 - x^2 - y^2},
\qquad x^2 + y^2 < R^2 ,
\]
regular em todo o seu domínio (\omterm{def:b2:metric:topology}{aberto}) --- inclusive no
polo que a carta esférica perdia. Perto de um ponto do equador
como $(R, 0, 0)$, o mesmo teorema resolve em $x$
($\frac{\partial f}{\partial x} = 2R \neq 0$). A regra
prática: uma superfície de nível é um gráfico sobre o plano coordenado
\emph{ortogonal à maior componente do gradiente} e,
cobrindo a esfera com seis dessas cartas gráfico,
verifica-se sua suavidade em toda parte sem trigonometria
alguma.
\end{example}

\section{Plano tangente e normal}

\begin{definition}[Plano tangente]\label{def:b2:surfaces:tangent}
Seja $\sigma$ regular em $(u_0, v_0)$, $M_0 = \sigma(u_0, v_0)$.
O \emph{plano tangente}\index{plano tangente} $T_{M_0}S$ é o plano por $M_0$
dirigido por $\operatorname{Vect}(\sigma_u, \sigma_v)$ (derivadas parciais em
$(u_0, v_0)$). A \emph{normal unitária}\index{normal unitária} é
\[
n(u_0, v_0)
= \frac{\sigma_u \wedge \sigma_v}
       {\norm{\sigma_u \wedge \sigma_v}} .
\]
\end{definition}

\begin{proposition}[Vetores tangentes são vetores
velocidade]\label{prop:b2:surfaces:velocity}
A direção de $T_{M_0}S$ é exatamente o conjunto dos vetores
$\gamma'(0)$, em que $\gamma = \sigma \circ c$ percorre as
curvas $\mathcal{C}^1$ desenhadas sobre a superfície passando por $M_0$ (isto é,
$c \colon (-\varepsilon, \varepsilon) \to U$ é $\mathcal{C}^1$
com $c(0) = (u_0, v_0)$).
\end{proposition}

\begin{proof}
Se $c(t) = (u(t), v(t))$, a regra da cadeia
(\cref{ch:b2:diffcalc}) dá
\[
\gamma'(0) = u'(0)\,\sigma_u + v'(0)\,\sigma_v
\in \operatorname{Vect}(\sigma_u, \sigma_v).
\]
Reciprocamente, o vetor $a\sigma_u + b\sigma_v$ é atingido pela
curva $c(t) = (u_0 + at,\ v_0 + bt)$, que permanece no \omterm{def:b2:metric:topology}{aberto}
$U$ para $\abs t$ pequeno.
\end{proof}

\begin{example}[O plano tangente do helicoide]
\label{ex:b2:surfaces:helicoidtangent}
Para o helicoide $\sigma(u, v) = (v\cos u,\ v\sin u,\ au)$,
no ponto $\sigma(0, 1) = (1, 0, 0)$:
\[
\sigma_u = (0,\ 1,\ a), \qquad \sigma_v = (1,\ 0,\ 0),
\qquad
\sigma_u \wedge \sigma_v = (0,\ a,\ -1),
\]
de modo que o \omterm{def:b2:surfaces:tangent}{plano tangente} é $a\,y = z$. Ele contém toda a
geratriz horizontal $\{(t, 0, 0)\}$ (direção $\sigma_v$): como
no cone do~\cref{exo:b2:surfaces:1}, uma superfície regrada por
retas tem cada geratriz contida no \omterm{def:b2:surfaces:tangent}{plano
tangente} ao longo dela. A outra direção tangente $\sigma_u$ é a
velocidade da hélice $u \mapsto \sigma(u, 1)$: uma carta,
duas curvas desenhadas, e o \omterm{def:b2:surfaces:tangent}{plano tangente} inteiro é \omterm{def:b2:structures:generated}{gerado} ---
a~\cref{prop:b2:surfaces:velocity} em ação.
\end{example}

\begin{proposition}[Normal de uma superfície de nível]\label{prop:b2:surfaces:gradient}
Seja $S = \{f = c\}$ com $f$ de classe $\mathcal{C}^1$ e $\nabla
f(M_0) \neq 0$. Então o \omterm{def:b2:surfaces:tangent}{plano tangente} de $S$ em $M_0$ é o
plano por $M_0$ ortogonal a $\nabla f(M_0)$:
\[
T_{M_0}S :\quad
\langle \nabla f(M_0),\ M - M_0\rangle = 0 .
\]
\end{proposition}

\begin{proof}
Para qualquer curva $\gamma$ desenhada sobre $S$ passando por $M_0$, $f(\gamma(t)) =
c$ identicamente, logo a regra da cadeia dá $\langle
\nabla f(M_0), \gamma'(0)\rangle = 0$: todos os vetores velocidade são
ortogonais ao gradiente, de modo que a direção tangente está contida
no plano $\nabla f(M_0)^\perp$. Ambos são subespaços de
dimensão $2$ --- a direção tangente porque $S$ é localmente um
gráfico regular (\cref{ex:b2:surfaces:standard}) e o complemento
ortogonal porque $\nabla f(M_0) \neq 0$ --- logo são
iguais.
\end{proof}

\begin{example}\label{ex:b2:surfaces:spheretangent}
Para a esfera $x^2 + y^2 + z^2 = R^2$: $\nabla f = 2(x, y, z)$,
de modo que o \omterm{def:b2:surfaces:tangent}{plano tangente} em $M_0$ é ortogonal ao raio
$\vect{OM_0}$ --- o fato clássico de que raio e \omterm{def:b2:surfaces:tangent}{plano tangente}
são perpendiculares, com equação $\langle M_0, M\rangle = R^2$.
\end{example}

\begin{example}[Plano tangente de um gráfico]\label{ex:b2:surfaces:graphtangent}
Para $z = f(x, y)$ em $(x_0, y_0)$: aplicando
a~\cref{prop:b2:surfaces:gradient} a $F(x,y,z) = f(x,y) - z$,
\[
z = f(x_0, y_0)
+ f_x(x_0, y_0)(x - x_0)
+ f_y(x_0, y_0)(y - y_0) ,
\]
a parte \omterm{def:b2:affine:subspace}{afim} do desenvolvimento de Taylor de primeira ordem --- o
\omterm{def:b2:surfaces:tangent}{plano tangente} é o gráfico da \omterm{def:b2:diffcalc:differential}{diferencial}, como deve ser.
\end{example}

\begin{example}[O ponto mais próximo de uma
superfície]\label{ex:b2:surfaces:closest}
Qual ponto do paraboloide $z = x^2 + y^2$ está mais próximo de
$P = (0, 0, 1)$? Minimize a distância ao quadrado ao longo da
superfície: com $\rho^2 = x^2 + y^2$,
\[
g(\rho^2) = \rho^2 + (\rho^2 - 1)^2,
\qquad
g'(\rho^2) = 1 + 2(\rho^2 - 1) = 0
\iff \rho^2 = \tfrac12 ,
\]
o que dá o círculo dos pontos de altura $z = \frac12$ e
distância $\sqrt{\tfrac12 + \tfrac14} = \frac{\sqrt3}2$. A
assinatura geométrica da minimalidade: num tal ponto $M$, o
vetor $\vect{MP}$ tem de ser \emph{normal} à superfície ---
caso contrário, deslizar ao longo de uma curva desenhada com velocidade tendo uma
componente na direção de $P$ diminuiria a distância. Confira:
$\nabla(z - x^2 - y^2) = (-2x, -2y, 1)$ em $M = (x, y,
\tfrac12)$, enquanto $\vect{MP} = (-x, -y, \tfrac12) =
\tfrac12(-2x, -2y, 1)$: paralelos, como previsto. A
condição de primeira ordem ``pé da perpendicular'' é a
mesma que governará os extremos sobre conjuntos de nível no
\cref{exo:b2:surfaces:12}.
\end{example}

\begin{example}[Planos tangentes de quádricas: a
regra da polarização]\label{ex:b2:surfaces:polarization}
Sejam $S : xy + yz + zx = 1$ e $M_0 = (1, 1, 0) \in S$. Aqui
$\nabla f = (y + z,\ x + z,\ x + y)$, logo $\nabla f(M_0) = (1,
1, 2)$ e o \omterm{def:b2:surfaces:tangent}{plano tangente} é
\[
(x - 1) + (y - 1) + 2z = 0,
\qquad\text{i.e.}\qquad x + y + 2z = 2 .
\]
A mesma resposta vem da \emph{regra de \omterm{def:b2:quadratic:def}{polarização}} que
generaliza \cref{ex:b2:surfaces:spheretangent} e
\cref{exo:b2:surfaces:2}: na equação da \omterm{pb:b2:surfaces:1}{quádrica}, substitua
$x^2$ por $x_0x$ e cada produto $xy$ por $\frac{x_0y +
y_0x}2$ (e ciclicamente):
\[
\frac{x_0y + y_0x}2 + \frac{y_0z + z_0y}2 + \frac{z_0x +
x_0z}2 = 1
\;\xrightarrow{\ M_0 = (1,1,0)\ }\;
\frac{x + y}2 + \frac z2 + \frac z2 = 1 ,
\]
que é de novo $x + y + 2z = 2$. A regra funciona porque
$\nabla$ de uma \omterm{def:b2:quadratic:def}{forma quadrática} é a forma bilinear associada
avaliada contra o ponto base --- tangência a uma \omterm{pb:b2:surfaces:1}{quádrica} é
\omterm{def:b2:quadratic:def}{polarização}, mais uma face do~\cref{ch:b2:quadratic}.
\end{example}

\section{A primeira forma fundamental}

\begin{definition}[Primeira forma fundamental]\label{def:b2:surfaces:fff}
Seja $\sigma \colon U \to \R^3$ uma superfície $\mathcal{C}^1$
regular. Sua \emph{primeira forma fundamental}\index{primeira forma
fundamental} em $(u,v)$ é a \omterm{def:b2:quadratic:def}{forma quadrática}
positiva definida em $\R^2$
\[
I(h, k) = \norm{h\,\sigma_u + k\,\sigma_v}^2
= E\,h^2 + 2F\,hk + G\,k^2,
\]
em que
\[
E = \norm{\sigma_u}^2,
\qquad
F = \langle \sigma_u, \sigma_v\rangle,
\qquad
G = \norm{\sigma_v}^2 .
\]
\end{definition}

\begin{remark}
$I$ é a restrição do produto escalar euclidiano ambiente ao
\omterm{def:b2:surfaces:tangent}{plano tangente}, lida na base $(\sigma_u, \sigma_v)$: ela é
positiva definida precisamente porque $\sigma_u, \sigma_v$ são
independentes (\cref{ch:b2:quadratic}). Toda quantidade métrica sobre a
superfície --- \omterm{def:b2:curves:length}{comprimentos} de curvas desenhadas, ângulos entre elas, \omterm{def:b2:surfaces:area}{áreas} ---
é calculada só a partir de $E, F, G$. Duas superfícies com a mesma $E,
F, G$ em parâmetros adequados são \emph{isométricas} mesmo que estejam
dispostas de modos diferentes no espaço: esse é o ponto de partida da geometria
intrínseca.
\end{remark}

\begin{example}[Ângulos entre as curvas
coordenadas]\label{ex:b2:surfaces:coordangle}
A \omterm{def:b2:surfaces:fff}{primeira forma fundamental} mede também ângulos: as
curvas coordenadas $u \mapsto \sigma(u, v_0)$ e $v \mapsto
\sigma(u_0, v)$ se encontram sob o ângulo $\theta$ com
\[
\cos\theta = \frac{\langle\sigma_u,
\sigma_v\rangle}{\norm{\sigma_u}\,\norm{\sigma_v}}
= \frac{F}{\sqrt{EG}} :
\]
o único coeficiente $F$ decide a ortogonalidade da
rede de parâmetros. Para a carta da esfera e para o helicoide, $F =
0$: meridianos cortam paralelos, e hélices cortam as geratrizes
horizontais, em ângulo reto --- e é por isso que seus integrandos de
\omterm{def:b2:surfaces:area}{área} colapsaram em $\sqrt{EG}$. Para uma carta gráfico,
$F = f_xf_y$ se anula apenas onde uma derivada parcial se anula:
a rede coordenada de um gráfico inclinado \emph{não} é
ortogonal, embora a rede em $(x, y)$ lá embaixo o seja. Quando os
cálculos sobre uma superfície ficam pesados, o primeiro movimento é
procurar uma carta com $F = 0$.
\end{example}

\begin{example}[A carta da sela]\label{ex:b2:surfaces:saddlechart}
Para a sela $z = xy$ com carta $\sigma(u, v) = (u, v,
uv)$:
\[
\sigma_u = (1, 0, v), \qquad \sigma_v = (0, 1, u),
\qquad
E = 1 + v^2, \quad F = uv, \quad G = 1 + u^2 ,
\]
e $EG - F^2 = 1 + u^2 + v^2 > 0$: regular em toda parte. As
duas curvas coordenadas por um ponto são \emph{retas}
de $\R^3$ (fixe $u$ ou fixe $v$: as geratrizes da
sela duplamente regrada), e no entanto $F \neq 0$ fora dos eixos: as geratrizes
por um ponto genérico não são ortogonais. Ambas as geratrizes estão
no \omterm{def:b2:surfaces:tangent}{plano tangente}, que elas geram --- de modo que o \omterm{def:b2:surfaces:tangent}{plano
tangente} corta a superfície ao longo de duas retas inteiras, o extremo
oposto da esfera, cujos \omterm{def:b2:surfaces:tangent}{planos tangentes} tocam num único
ponto. O sinal do ``contato de segunda ordem'' entre
uma superfície e seus \omterm{def:b2:surfaces:tangent}{planos tangentes} é uma história de \omterm{thm:b2:curves:frenet2d}{curvatura},
retomada no volume do terceiro ano de graduação.
\end{example}

\begin{proposition}[Comprimento de uma curva desenhada sobre uma
superfície]\label{prop:b2:surfaces:curvelength}
Se $\gamma(t) = \sigma(u(t), v(t))$, $t \in [a, b]$, é
$\mathcal{C}^1$, então
\[
L(\gamma) = \int_a^b
\sqrt{E\,u'^2 + 2F\,u'v' + G\,v'^2}\;\dd t ,
\]
com $E, F, G$ avaliados em $(u(t), v(t))$.
\end{proposition}

\begin{proof}
$\gamma' = u'\sigma_u + v'\sigma_v$ pela regra da cadeia, logo
$\norm{\gamma'}^2 = I(u', v') = Eu'^2 + 2Fu'v' + Gv'^2$; integre
$\norm{\gamma'}$ (\cref{def:b2:curves:length}).
\end{proof}

\begin{example}[Por que os aviões de linha voam sobre o polo]\label{ex:b2:surfaces:polarroute}
Dois aeroportos estão na latitude $\varphi_0$ e em longitudes
opostas: $A = \sigma(0, \varphi_0)$ e $B = \sigma(\pi,
\varphi_0)$ na esfera de raio $R$. Ao longo do paralelo
($\varphi \equiv \varphi_0$), o \omterm{def:b2:curves:length}{comprimento} é
$\int_0^\pi R\cos\varphi_0\,\dd\theta = \pi R\cos\varphi_0$.
Ao longo da rota sobre o polo (subindo o meridiano $\theta = 0$,
descendo o meridiano $\theta = \pi$), ele vale $2R(\frac\pi2 -
\varphi_0)$. Na latitude $\varphi_0 = \frac\pi3$ (sessenta
graus): rota do paralelo $\pi R/2 \approx 1.571\,R$, rota
polar $\pi R/3 \approx 1.047\,R$ --- um terço mais curta. De fato,
$\pi\cos\varphi_0 \geq \pi - 2\varphi_0$ em
$\intcc0{\pi/2}$ (a função $\pi\cos\varphi - \pi +
2\varphi$ se anula nas duas extremidades e sua derivada $2 -
\pi\sin\varphi$ muda de sinal uma única vez, logo ela é primeiro crescente
e depois decrescente, portanto não negativa): a rota polar nunca
perde. A \omterm{def:b2:surfaces:fff}{primeira forma fundamental} transformou uma questão de navegação
em duas integrais de uma linha; o~\cref{exo:b2:surfaces:6} leva
a ideia até uma verdadeira demonstração de minimalidade para os meridianos.
\end{example}

\begin{lemma}[Identidade de Lagrange]\label{lem:b2:surfaces:lagrange}
Para todos $a, b \in \R^3$:
$\norm{a \wedge b}^2 = \norm a^2 \norm b^2 - \langle a, b\rangle^2$.
Em particular
\[
\norm{\sigma_u \wedge \sigma_v} = \sqrt{EG - F^2}.
\]
\end{lemma}

\begin{proof}
Ambos os membros ficam inalterados se substituirmos $b$ por sua componente $b_\perp
= b - \frac{\langle a, b\rangle}{\norm a^2}a$ ortogonal a $a$
(para $a \neq 0$; o caso $a = 0$ é trivial): o membro da esquerda
porque $a \wedge a = 0$, o da direita desenvolvendo $\norm
{b_\perp}^2 = \norm b^2 - \frac{\langle a,b\rangle^2}{\norm a^2}$
e $\langle a, b_\perp\rangle = 0$. Basta portanto demonstrar a
identidade para $a, b$ ortogonais, em que ela se lê $\norm{a \wedge
b} = \norm a \norm b$: verdadeiro, pois para vetores ortogonais o
produto vetorial tem \omterm{def:b2:nvs:norm}{norma} $\norm a\norm b\,\abs{\sin\frac\pi2}$. A
fórmula exibida é o caso $a = \sigma_u$, $b = \sigma_v$.
\end{proof}

\begin{remark}
A identidade de Lagrange diz que $EG - F^2 = \det\begin{psmallmatrix}
E & F\\ F & G\end{psmallmatrix}$ é o \emph{\omterm{def:b2:linalg:det}{determinante}
de Gram} de $(\sigma_u, \sigma_v)$: a \omterm{def:b2:surfaces:area}{área} ao quadrado do
paralelogramo por eles \omterm{def:b2:structures:generated}{gerado}. Regularidade, positividade
definida da \omterm{def:b2:surfaces:fff}{primeira forma fundamental} e positividade do
\omterm{def:b2:linalg:det}{determinante} de Gram são três formulações de uma só condição ---
e é por isso que o integrando de \omterm{def:b2:surfaces:area}{área} abaixo nunca se anula numa
carta regular.
\end{remark}

\begin{definition}[Área]\label{def:b2:surfaces:area}
Sejam $\sigma \colon U \to \R^3$ uma superfície $\mathcal{C}^1$ regular e injetiva
e $K \subseteq U$ um domínio \omterm{def:b2:metric:compact}{compacto} sobre o qual as integrais duplas
façam sentido (\cref{ch:b2:multint}). A
\emph{área}\index{área!de uma superfície} do pedaço $\sigma(K)$ é
\[
\mathcal{A}
= \iint_K \norm{\sigma_u \wedge \sigma_v}\,\dd u\,\dd v
= \iint_K \sqrt{EG - F^2}\;\dd u\,\dd v .
\]
\end{definition}

\begin{remark}[Por que essa fórmula]
O retângulo $[u, u + \dd u] \times [v, v + \dd v]$ é aplicado, em
primeira ordem, sobre o paralelogramo \omterm{def:b2:structures:generated}{gerado} por $\sigma_u\,\dd u$
e $\sigma_v\,\dd v$, cuja \omterm{def:b2:surfaces:area}{área} vale $\norm{\sigma_u \wedge
\sigma_v}\,\dd u\,\dd v$: a definição integra o fator local
de distorção de \omterm{def:b2:surfaces:area}{área}, exatamente como o \omterm{def:b2:curves:length}{comprimento de arco} integra a
velocidade local. A compatibilidade com \omterm{def:b2:curves:reparam}{mudanças de parâmetro} é o
\cref{exo:b2:surfaces:8}; a compatibilidade com a fórmula de mudança de
variáveis para integrais duplas é discutida no
\cref{ch:b2:multint}.
\end{remark}

\begin{example}[Dois gráficos diferentes, uma só
área]\label{ex:b2:surfaces:samearea}
Sobre o disco unitário, compare a tigela $z = \frac12(x^2 + y^2)$
e a sela $z = xy$. Seus integrandos de \omterm{def:b2:surfaces:area}{área}
(\cref{exo:b2:surfaces:5}) são
\[
\sqrt{1 + x^2 + y^2}
\qquad\text{e}\qquad
\sqrt{1 + y^2 + x^2} :
\]
idênticos. As duas superfícies --- uma se curvando do mesmo modo em
todas as direções, a outra em forma de sela --- têm exatamente
as mesmas \omterm{def:b2:surfaces:area}{áreas} sobre todo domínio, $\frac{2\pi}3(2\sqrt2 - 1)$
sobre o disco unitário. O elemento de \omterm{def:b2:surfaces:area}{área} só enxerga o
\emph{\omterm{def:b2:curves:length}{comprimento}} do gradiente, não a disposição da
flexão; distinguir a tigela da sela exige dados de
segunda ordem (a estrutura de sinais exibida no
\cref{fig:b2:surfaces:saddle}), que nenhuma medida de \omterm{def:b2:surfaces:area}{área}
detecta. \omterm{def:b2:surfaces:fff}{Primeira forma fundamental}: métrica, cega
à forma; a segunda forma, que vê a forma, pertence ao terceiro ano.
\end{example}

\begin{example}[Área da esfera]\label{ex:b2:surfaces:spherearea}
Para a carta esférica de \cref{ex:b2:surfaces:standard}:
\[
\sigma_\theta = R(-\sin\theta\cos\varphi,\
\cos\theta\cos\varphi,\ 0),
\qquad
\sigma_\varphi = R(-\cos\theta\sin\varphi,\
-\sin\theta\sin\varphi,\ \cos\varphi),
\]
logo $E = R^2\cos^2\varphi$, $F = 0$, $G = R^2$ e $\sqrt{EG - F^2}
= R^2\cos\varphi$. Portanto
\[
\mathcal{A}
= \int_{-\pi/2}^{\pi/2}\!\!\int_0^{2\pi}
R^2\cos\varphi\;\dd\theta\,\dd\varphi
= 2\pi R^2\,\bigl[\sin\varphi\bigr]_{-\pi/2}^{\pi/2}
= \boxed{4\pi R^2} .
\]
\end{example}

\begin{example}[O cone, conferido contra a
fórmula da escola]\label{ex:b2:surfaces:conearea}
Para o cone $z = \sqrt{x^2 + y^2}$ sobre a coroa $a \leq
\rho \leq b$, a fórmula do gráfico do
\cref{exo:b2:surfaces:5} dá $1 + f_x^2 + f_y^2 = 2$
(calcule $f_x = x/\rho$, $f_y = y/\rho$), logo
\[
\mathcal A = \sqrt2\,\pi\,(b^2 - a^2) .
\]
Verificação de coerência com a fórmula da geratriz $\pi\rho\ell$
do~\cref{ex:b2:surfaces:revolution}: os cones \omterm{def:b2:metric:complete}{completos} de raios de
base $b$ e $a$ têm \omterm{def:b2:surfaces:area}{áreas} laterais $\pi b\cdot b\sqrt2$ e
$\pi a\cdot a\sqrt2$, cuja diferença é exatamente
$\sqrt2\pi(b^2 - a^2)$. Duas cartas, duas fórmulas, uma só \omterm{def:b2:surfaces:area}{área}
--- a invariância demonstrada no~\cref{exo:b2:surfaces:8}, vista
em pleno funcionamento.
\end{example}

\begin{remark}[Verificações de bom senso para áreas]
Três verificações instantâneas pegam a maioria dos erros num cálculo
de \omterm{def:b2:surfaces:area}{área}. \emph{Escala}: dilatar uma superfície por $\lambda$
multiplica $E, F, G$ por $\lambda^2$ e a \omterm{def:b2:surfaces:area}{área} por
$\lambda^2$ --- uma resposta cuja dependência em $R$ não seja
quadrática (como $4\pi R^2$) está errada. \emph{Positividade do
elemento}: $\sqrt{EG - F^2}$ tem de ser estritamente positivo no
interior da carta; um valor nulo sinaliza uma degenerescência da carta,
a ser excisada como nos polos da esfera. \emph{Simetria}: um
cálculo sobre um pedaço \omterm{def:b2:quadratic:adjoint}{simétrico} tem de ser compatível com a
soma de suas partes congruentes --- é bom que o hemisfério
dê $2\pi R^2$.
\end{remark}

\begin{example}[Superfície de revolução]\label{ex:b2:surfaces:revolution}
Rode a curva $z \mapsto (r(z), 0, z)$, $r > 0$ de classe
$\mathcal{C}^1$, em torno do eixo $z$:
\[
\sigma(\theta, z) = (r(z)\cos\theta,\ r(z)\sin\theta,\ z).
\]
Então $E = r(z)^2$, $F = 0$, $G = 1 + r'(z)^2$, logo
\[
\mathcal{A} = \int_a^b\!\!\int_0^{2\pi}
r(z)\sqrt{1 + r'(z)^2}\;\dd\theta\,\dd z
= 2\pi\int_a^b r(z)\sqrt{1 + r'(z)^2}\,\dd z ,
\]
a fórmula clássica (circunferência $2\pi r$ vezes o elemento de
\omterm{def:b2:curves:length}{comprimento} da geratriz). Para o cone $r(z) = kz$, $z \in [0, h]$: $\mathcal
A = 2\pi k\sqrt{1 + k^2}\,\frac{h^2}{2} = \pi \rho \ell$ com
$\rho = kh$ o raio da base e $\ell = h\sqrt{1 + k^2}$ a geratriz
--- a fórmula da escola, agora demonstrada e não admitida.
\end{example}

\begin{example}[O catenoide]\label{ex:b2:surfaces:catenoid}
Rode a catenária $r(z) = \cosh z$, $z \in \intcc{-1}{1}$,
em torno de seu eixo: o \emph{catenoide} resultante tem, pela
fórmula de revolução e pelo $1 + \sinh^2 z = \cosh^2 z$,
\[
\mathcal A = 2\pi\int_{-1}^1\cosh z\,\sqrt{1 + \sinh^2z}\,
\dd z = 2\pi\int_{-1}^1\cosh^2z\,\dd z
= \pi\bigl[z + \sinh z\cosh z\bigr]_{-1}^{1},
\]
isto é
\[
\mathcal A = 2\pi + \pi\sinh 2 \approx 17.68 .
\]
O fator da geratriz $\sqrt{1 + r'^2}$ fundiu-se com o raio
num quadrado perfeito --- a mesma identidade que tornou elementar o
\omterm{def:b2:curves:length}{comprimento de arco} da catenária no capítulo das curvas. Isso
não é um acaso algébrico: entre todas as superfícies de revolução
que se apoiam nos dois círculos de bordo, o catenoide minimiza a
\omterm{def:b2:surfaces:area}{área} (é a forma de uma película de sabão entre dois anéis), e
essa propriedade variacional é precisamente o que singulariza
$\cosh$; o volume do terceiro ano de graduação a demonstra com o cálculo das
variações.
\end{example}

\begin{figure}[ht]
\centering
\begin{tikzpicture}[scale=1.05]
  % A saddle surface z = x^2 - y^2 drawn as a wire mesh in oblique projection
  % projection: (x,y,z) -> (x + 0.45*y, z + 0.35*y)
  \begin{scope}
    % u-lines (x varies, y fixed)
    \foreach \yy in {-1,-0.5,0,0.5,1} {
      \draw[ocDarkBlue!70, thick, domain=-1:1, smooth, samples=25]
        plot ({\x + 0.45*\yy}, {\x*\x - \yy*\yy + 0.35*\yy});
    }
    % v-lines (y varies, x fixed)
    \foreach \xx in {-1,-0.5,0,0.5,1} {
      \draw[ocMint!80!black, thick, domain=-1:1, smooth, samples=25]
        plot ({\xx + 0.45*\x}, {\xx*\xx - \x*\x + 0.35*\x});
    }
    % tangent plane at origin: z = 0 -> parallelogram
    \draw[ocRed, dashed]
      (-0.8 - 0.36, -0.28) -- (0.8 - 0.36, -0.28)
      -- (0.8 + 0.36, 0.28) -- (-0.8 + 0.36, 0.28) -- cycle;
    \fill (0,0) circle (0.045) node[below right, font=\small] {$M_0$};
    % normal vector at origin
    \draw[->, thick] (0,0) -- (0,0.9)
      node[right, font=\small] {$n$};
  \end{scope}
\end{tikzpicture}
\caption{A sela $z = x^2 - y^2$ perto da origem, com suas
curvas coordenadas (curvas $u$ em azul, curvas $v$ em verde), o
\omterm{def:b2:surfaces:tangent}{plano tangente} em $M_0 = (0,0,0)$ (tracejado) e a normal unitária $n$.
A superfície atravessa seu \omterm{def:b2:surfaces:tangent}{plano tangente} --- o análogo
bidimensional de uma inflexão.}
\label{fig:b2:surfaces:saddle}
\end{figure}

\begin{remark}[Armadilhas comuns]
(i) \emph{Singularidades da carta não são singularidades da superfície}:
a carta esférica degenera nos polos
($\cos\varphi = 0$), mas a esfera é perfeitamente suave ali
--- outra carta (troque os papéis dos eixos) é regular
nos polos. Antes de declarar um ponto singular, tente uma segunda
parametrização. (ii) \emph{A regularidade de $\sigma$ diz respeito
à parametrização, não à imagem}: $\sigma(u, v) = (u^3, v,
0)$ não é regular ao longo de $u = 0$ embora sua imagem seja um
plano. (iii) A fórmula da \omterm{def:b2:surfaces:area}{área} exige $\sigma$
\emph{injetiva} em $K$: uma carta que cobre um pedaço duas vezes
o conta duas vezes ($\theta$ percorrendo $\intcc0{4\pi}$
duplica a \omterm{def:b2:surfaces:area}{área} da esfera). (iv) A normal unitária é definida
a menos de sinal pela superfície, mas é \emph{escolhida} pela carta
(\omterm{def:b2:structures:generated}{ordem} de $u, v$); enunciados que envolvem orientação têm de fixar
essa escolha. (v) Por fim, $EG - F^2 > 0$ não é uma hipótese
suplementar: é exatamente a regularidade, pela identidade de
Lagrange --- se ela se anula em algum lugar, o problema é a
carta, e nenhuma fórmula de \omterm{def:b2:surfaces:area}{área} ou de \omterm{def:b2:surfaces:tangent}{plano tangente} se aplica ali.
\end{remark}

\begin{remark}[Perspectivas dentro deste volume]
Elos para a frente a partir daqui. O elemento de \omterm{def:b2:surfaces:area}{área} $\sqrt{EG -
F^2}\,\dd u\,\dd v$ é um jacobiano bidimensional
disfarçado, e o~\cref{ch:b2:multint} torna a analogia exata
com o teorema de mudança de variáveis --- as integrais de
superfície de lá são as \omterm{def:b2:surfaces:area}{áreas} deste capítulo com um integrando a
bordo. A \omterm{def:b2:surfaces:fff}{primeira forma fundamental} é um campo de
\omterm{def:b2:quadratic:def}{formas quadráticas} positivas, tratado \omterm{def:b2:funcseq:def}{pontualmente} pelas ferramentas do
\cref{ch:b2:quadratic}, cujo teorema espectral também move a
classificação das \omterm{pb:b2:surfaces:1}{quádricas} do problema de fim de semana. E a
reta normal governa problemas de extremo sobre conjuntos de restrição
(\cref{ex:b2:surfaces:closest}), o germe geométrico do
método dos multiplicadores de Lagrange esboçado com o
\cref{thm:b2:diffcalc:extrema}.
\end{remark}

\begin{omfigure}
\begin{tikzpicture}[scale=0.8]
  \draw[omDef, thick] (0,0) ellipse (1.5 and 0.95);
  \draw[omThm] (0,0) ellipse (1.5 and 0.35);
  \node[below, font=\small] at (0,-1.25) {elipsoide};
\end{tikzpicture}
\qquad
\begin{tikzpicture}[scale=0.8]
  \draw[omDef, thick] plot[domain=-1.1:1.1, samples=30]
    ({0.6*(exp(\x)+exp(-\x))/2}, {\x});
  \draw[omDef, thick] plot[domain=-1.1:1.1, samples=30]
    ({-0.6*(exp(\x)+exp(-\x))/2}, {\x});
  \draw[omThm] (0,0) ellipse (0.6 and 0.16);
  \draw[omDef] (0,1.1) ellipse (1.0 and 0.22);
  \draw[omDef] (0,-1.1) ellipse (1.0 and 0.22);
  \node[below, font=\small] at (0,-1.55)
    {hiperboloide de uma folha};
\end{tikzpicture}
\qquad
\begin{tikzpicture}[scale=0.8]
  \draw[omDef, thick] plot[domain=-1.05:1.05, samples=30]
    ({\x}, {1.1*\x*\x - 1.1});
  \draw[omThm] (0,0.055) ellipse (1.0 and 0.22);
  \node[below, font=\small] at (0,-1.45)
    {paraboloide elíptico};
\end{tikzpicture}
\qquad
\begin{tikzpicture}[scale=0.8]
  \draw[omDef, thick] plot[domain=-1:1, samples=30]
    ({\x}, {0.55*\x*\x});
  \draw[omThm, thick] plot[domain=-1:1, samples=30]
    ({0.45*\x}, {-0.55*\x*\x + 0.18*\x});
  \node[below, font=\small] at (0,-1.15)
    {paraboloide hiperbólico};
\end{tikzpicture}

{\small Quatro das nove superfícies \omterm{pb:b2:surfaces:1}{quádricas} classificadas no
problema de fim de semana, esboçadas por suas silhuetas e uma curva
de nível (vermelho): o elipsoide limitado, o hiperboloide de uma
folha duplamente regrado com sua cintura, a tigela do
paraboloide elíptico e a sela, cujas duas seções
parabólicas se dobram em sentidos opostos.}
\end{omfigure}

\begin{remark}[Onde isso é usado]
O elemento de \omterm{def:b2:surfaces:area}{área} $\norm{\sigma_u\wedge\sigma_v}\,\dd u\,\dd v$
é a medida das integrais de superfície do~\cref{ch:b2:multint}, em que
ele encontra a fórmula de Green; a \omterm{def:b2:surfaces:fff}{primeira forma fundamental} é o
protótipo de um campo de \omterm{def:b2:quadratic:def}{formas quadráticas}, estudado \omterm{def:b2:funcseq:def}{pontualmente}
com as ferramentas do~\cref{ch:b2:quadratic}; e o problema de fim de
semana deste capítulo classifica todas as superfícies \omterm{pb:b2:surfaces:1}{quádricas} com o
teorema espectral. O volume do terceiro ano de graduação volta às superfícies com
formas \omterm{def:b2:diffcalc:differential}{diferenciais} e o teorema da divergência, e a \omterm{thm:b2:curves:frenet2d}{curvatura}
intrínseca --- o que $E, F, G$ sabem sobre a flexão --- é a
porta de entrada da geometria diferencial propriamente dita.
\end{remark}

\section{Exercícios}

\begin{exercise}[$\star$]\label{exo:b2:surfaces:1}
Mostre que os \omterm{def:b2:surfaces:tangent}{planos tangentes} do cone $z = \sqrt{x^2 + y^2}$
(menos seu vértice) passam todos pelo vértice. \emph{(Parametrize por
$\sigma(\theta, r) = (r\cos\theta, r\sin\theta, r)$, $r > 0$.)}
\end{exercise}

\begin{exercise}[$\star$]\label{exo:b2:surfaces:2}
Determine o \omterm{def:b2:surfaces:tangent}{plano tangente} do elipsoide $\frac{x^2}{a^2} +
\frac{y^2}{b^2} + \frac{z^2}{c^2} = 1$ num ponto $(x_0, y_0,
z_0)$ da superfície.
\end{exercise}

\begin{exercise}[$\star$]\label{exo:b2:surfaces:3}
Calcule $E, F, G$ para o helicoide $\sigma(u, v) = (v\cos u,\
v\sin u,\ au)$, $a > 0$, e a \omterm{def:b2:surfaces:area}{área} do pedaço $0 \leq u \leq
2\pi$, $0 \leq v \leq 1$, como uma integral (avalie-a usando
$\int\sqrt{v^2 + a^2}\,\dd v = \frac12\bigl(v\sqrt{v^2+a^2} +
a^2\ln(v + \sqrt{v^2 + a^2})\bigr) + C$).
\end{exercise}

\begin{exercise}[$\star\star$]\label{exo:b2:surfaces:4}
(Toro) Parametrize o toro obtido girando o círculo de
centro $(R, 0, 0)$ e raio $r < R$ no plano $xz$ em torno do
eixo $z$:
\[
\sigma(\theta, \psi) = \bigl((R + r\cos\psi)\cos\theta,\
(R + r\cos\psi)\sin\theta,\ r\sin\psi\bigr).
\]
Calcule $E, F, G$, verifique a regularidade e mostre que a \omterm{def:b2:surfaces:area}{área} vale
$4\pi^2 R r$ (Pappus: circunferência média $2\pi R$ vezes o
\omterm{def:b2:curves:length}{comprimento} do círculo $2\pi r$).
\end{exercise}

\begin{exercise}[$\star\star$]\label{exo:b2:surfaces:5}
Mostre que a \omterm{def:b2:surfaces:area}{área} do gráfico de $f \in \mathcal{C}^1(K)$ é
$\iint_K \sqrt{1 + f_x^2 + f_y^2}\,\dd x\,\dd y$ e calcule-a
para o pedaço de paraboloide $z = \frac12(x^2 + y^2)$ sobre o disco
$x^2 + y^2 \leq 1$ (coordenadas polares,
\cref{ch:b2:multint}).
\end{exercise}

\begin{exercise}[$\star\star$]\label{exo:b2:surfaces:6}
Uma curva desenhada $\gamma(t) = \sigma(u(t), v(t))$ sobre a esfera de
raio $R$ (carta esférica) tem $u = \theta(t)$, $v =
\varphi(t)$. Escreva seu \omterm{def:b2:curves:length}{comprimento} como uma integral em $\theta, \varphi$
e demonstre que, entre as curvas que ligam dois pontos de um mesmo
meridiano $\theta = \theta_0$, o arco de meridiano é a mais curta.
\emph{(Minore o integrando por $R\abs{\varphi'}$.)}
\end{exercise}

\begin{exercise}[$\star\star\star$]\label{exo:b2:surfaces:7}
(Retas normais de uma esfera) Seja $S$ uma superfície de nível regular
$\{f = c\}$, \omterm{def:b2:metric:connected}{conexa}, cujas retas normais passam todas por um
ponto fixo $\Omega$. Mostre que $S$ está contida numa esfera
centrada em $\Omega$. \emph{(Mostre que $\norm{M - \Omega}^2$ tem
derivada nula ao longo de toda curva desenhada sobre $S$.)}
\end{exercise}

\begin{exercise}[$\star\star\star$]\label{exo:b2:surfaces:8}
(A \omterm{def:b2:surfaces:area}{área} é geométrica) Seja $\Phi \colon U' \to U$ um difeomorfismo
$\mathcal{C}^1$ entre \omterm{def:b2:metric:topology}{abertos} de $\R^2$ e $\tilde\sigma =
\sigma \circ \Phi$. Mostre que
\[
\norm{\tilde\sigma_{u'} \wedge \tilde\sigma_{v'}}
= \abs{\det J_\Phi}\,
\norm{(\sigma_u \wedge \sigma_v)\circ\Phi} ,
\]
e deduza, usando a fórmula de mudança de variáveis do
\cref{ch:b2:multint}, que a \omterm{def:b2:surfaces:area}{área} de
\cref{def:b2:surfaces:area} não depende da parametrização regular
escolhida.
\end{exercise}

\begin{exercise}[$\star$]\label{exo:b2:surfaces:9}
(Teorema da caixa de chapéu de Arquimedes) Na esfera de raio $R$, a
\emph{zona} entre as latitudes com $z_1 \leq z \leq z_2$
($-R \leq z_1 < z_2 \leq R$) tem \omterm{def:b2:surfaces:area}{área} $2\pi R\,(z_2 - z_1)$:
demonstre-o com a carta esférica e conclua que a \omterm{def:b2:surfaces:area}{área} de uma zona
depende apenas de sua altura --- fatiar uma laranja em fatias de
espessura igual dá quantidades iguais de casca.
\end{exercise}

\begin{exercise}[$\star\star$]\label{exo:b2:surfaces:10}
Mostre que toda reta normal de uma superfície de revolução
$\sigma(\theta, z) = (r(z)\cos\theta,\ r(z)\sin\theta,\ z)$
($r > 0$ de classe $\mathcal C^1$) encontra o eixo de
revolução, e localize o ponto de interseção.
\end{exercise}

\begin{exercise}[$\star\star$]\label{exo:b2:surfaces:11}
(Desenrolando o cilindro) A carta $\sigma(u, v) = (\cos u,\
\sin u,\ v)$ do cilindro unitário tem $E = G = 1$, $F = 0$:
verifique isso e explique por que toda curva desenhada $t \mapsto
\sigma(u(t), v(t))$ tem o mesmo \omterm{def:b2:curves:length}{comprimento} que a curva plana $t
\mapsto (u(t), v(t))$. Deduza que a hélice que vai de $(1, 0, 0)$
a $(1, 0, 2\pi c)$ dando uma volta tem \omterm{def:b2:curves:length}{comprimento}
$2\pi\sqrt{1 + c^2}$ e que nenhuma curva desenhada com os mesmos
extremos e uma volta completa é mais curta.
\end{exercise}

\begin{exercise}[$\star\star\star$]\label{exo:b2:surfaces:12}
Sejam $S = \{f = c\}$ uma superfície de nível regular \omterm{def:b2:metric:compact}{compacta} e
$M_0 \in S$ um ponto à distância máxima da origem. Mostre
que $\nabla f(M_0)$ é colinear com $\vect{OM_0}$ --- a
normal no ponto mais distante é radial. Aplique ao elipsoide
$\frac{x^2}{a^2} + \frac{y^2}{b^2} + \frac{z^2}{c^2} = 1$ ($a >
b > c > 0$): determine todos os pontos em que a normal é radial e
identifique os mais distantes.
\end{exercise}

\section{Problema: a classificação das quádricas de
$\R^3$}

\begin{problem}[{Problema de fim de semana --- toda superfície quádrica,
ordenada pelo teorema espectral}]\label{pb:b2:surfaces:1}
Uma \emph{quádrica}\index{quádrica} é o conjunto de zeros em $\R^3$ de um
polinômio de grau dois
\[
q(X) = X^{\mathsf T}\!AX + 2\,\langle b, X\rangle + c,
\qquad A \in \mathcal S_3(\R),\ A \neq 0,\ b \in \R^3,\
c \in \R .
\]
As superfícies das figuras deste capítulo --- esferas,
elipsoides, selas, cones, cilindros --- são todas \omterm{pb:b2:surfaces:1}{quádricas}.
Este problema as classifica \emph{todas}: o teorema espectral
(\cref{thm:b2:quadratic:spectral}) endireita a parte quadrática,
translações \omterm{def:b2:affine:subspace}{afins} (\cref{ch:b2:affine}) absorvem a
parte linear, e o que resta é uma lista curta e \omterm{def:b2:metric:complete}{completa} de
formas normais.

\medskip
\noindent\textbf{Parte I --- A máquina de redução.}
\begin{enumerate}
  \item Seja $X = PY + t$ com $P \in O(3)$ e $t \in \R^3$ (uma
        mudança rígida de coordenadas). Mostre que $q(PY + t) =
        Y^{\mathsf T}\!A'Y + 2\langle b', Y\rangle + c'$ com
        \[
        A' = P^{\mathsf T}\!AP, \qquad
        b' = P^{\mathsf T}(At + b), \qquad
        c' = q(t) .
        \]
        Deduza que o \omterm{def:b2:reduction:eigen}{espectro} de $A$ (logo seu posto e sua
        assinatura) é um invariante rígido da equação e
        explique por que a equação de uma dada \omterm{pb:b2:surfaces:1}{quádrica} só é
        determinada a menos de um fator escalar não nulo.
  \item Usando o teorema espectral, mostre que, após uma rotação,
        a equação se torna $\sum_i \lambda_i y_i^2 +
        2\sum_i\beta_iy_i + c = 0$ com $\lambda_1, \lambda_2,
        \lambda_3$ os \omterm{def:b2:reduction:eigen}{autovalores} de $A$.
  \item Para todo $i$ com $\lambda_i \neq 0$, absorva
        $\beta_iy_i$ por uma translação ($y_i \mapsto y_i -
        \beta_i/\lambda_i$). Escreva a equação reduzida quando
        $\operatorname{rank} A = r$: $\sum_{i\leq r}\lambda_i
        z_i^2 + 2\sum_{i > r}\beta_iz_i + c'' = 0$.
  \item Um \emph{centro} da \omterm{pb:b2:surfaces:1}{quádrica} de equação $q = 0$ é
        um ponto $\Omega$ com $q(2\Omega - X) = q(X)$ para todo
        $X$: a simetria \omterm{def:b2:funcseq:def}{pontual} em $\Omega$ preserva a
        equação, logo a superfície. Mostre que $q(2\Omega - X)
        - q(X) = -4\langle A\Omega + b, X\rangle + 4\langle
        A\Omega + b, \Omega\rangle$ e deduza: os centros
        são exatamente as soluções de $A\Omega = -b$; eles
        existem se e somente se $b \in \operatorname{im}A$, e o centro
        é único se e somente se $A$ é invertível.
\end{enumerate}

\medskip
\noindent\textbf{Parte II --- \omterm{pb:b2:surfaces:1}{Quádricas} centrais
($\operatorname{rank}A = 3$).} Aqui a equação reduzida é
$\lambda_1z_1^2 + \lambda_2z_2^2 + \lambda_3z_3^2 = \delta$.
\begin{enumerate}[resume]
  \item Multiplicando por $-1$ se necessário, suponha ao menos dois
        $\lambda_i > 0$. Enumere as possibilidades:
        assinatura $(3, 0)$ com $\delta > 0$, $= 0$, $< 0$,
        e assinatura $(2, 1)$ com $\delta > 0$, $= 0$, $<
        0$; nomeie os seis conjuntos resultantes (elipsoide, ponto,
        conjunto vazio, hiperboloide de uma folha, cone, hiperboloide
        de duas folhas) e ponha cada um em sua forma normal
        euclidiana ($\frac{x^2}{a^2} + \frac{y^2}{b^2} +
        \frac{z^2}{c^2} = 1$ etc.).
  \item Classifique $x^2 + y^2 + z^2 + 4xy + 4yz + 4zx = 1$:
        mostre que $A = 2J - I$ com $J$ a matriz de uns,
        calcule o \omterm{def:b2:reduction:eigen}{espectro} $\{5, -1, -1\}$ e identifique um
        hiperboloide de duas folhas de revolução em torno do eixo
        $\R(1,1,1)$.
  \item (Geratrizes) Para o hiperboloide de uma folha
        $\frac{x^2}{a^2} + \frac{y^2}{b^2} - \frac{z^2}{c^2}
        = 1$, fatore
        \[
        \Bigl(\frac xa - \frac zc\Bigr)
        \Bigl(\frac xa + \frac zc\Bigr)
        = \Bigl(1 - \frac yb\Bigr)\Bigl(1 + \frac yb\Bigr)
        \]
        e produza duas famílias a um parâmetro de retas
        contidas na superfície.
  \item Mostre que por todo ponto do hiperboloide de uma
        folha passa exatamente uma reta de cada família:
        a superfície é \emph{duplamente regrada}.
  \item O \emph{cone assintótico} do hiperboloide de uma
        folha é $C : \frac{x^2}{a^2} + \frac{y^2}{b^2}
        - \frac{z^2}{c^2} = 0$. Com $\rho^2 = \frac{x^2}{a^2}
        + \frac{y^2}{b^2}$, mostre que todo ponto do
        hiperboloide está à distância no máximo $c\bigl(\rho -
        \sqrt{\rho^2 - 1}\bigr) = \frac{c}{\rho +
        \sqrt{\rho^2 - 1}}$ de $C$, de modo que a superfície se cola ao
        seu cone no infinito. Quais são as seções do
        hiperboloide pelos planos $x = \pm a$?
\end{enumerate}

\medskip
\noindent\textbf{Parte III --- Posto $2$ e posto $1$:
paraboloides, cilindros, planos.}
\begin{enumerate}[resume]
  \item Suponha $\operatorname{rank}A = 2$, digamos $\lambda_1,
        \lambda_2 \neq 0 = \lambda_3$. Partindo da questão
        3, separe em dois casos conforme $\beta_3 \neq 0$
        (sem centro, pela questão 4) ou $\beta_3 = 0$ (uma reta
        de centros), e reduza a
        \[
        \lambda_1z_1^2 + \lambda_2z_2^2 + 2\beta_3z_3 = 0
        \qquad\text{ou}\qquad
        \lambda_1z_1^2 + \lambda_2z_2^2 + c'' = 0 :
        \]
        paraboloides elípticos/hiperbólicos no primeiro caso,
        cilindros sobre cônicas centrais (ou pares de
        planos concorrentes, uma reta, o conjunto vazio) no
        segundo.
  \item Mostre que a sela $z = xy$ é um paraboloide
        hiperbólico: rode de $\pi/4$ no plano $xy$ para
        chegar a $z = \tfrac12(u^2 - v^2)$, a superfície do
        \cref{fig:b2:surfaces:saddle} a menos de escala.
  \item Mostre que a sela $z = xy$ carrega as duas famílias de
        retas $\{x = x_0,\ z = x_0y\}$ e $\{y = y_0,\ z
        = xy_0\}$, com exatamente uma reta de cada uma passando por
        todo ponto: a segunda \omterm{pb:b2:surfaces:1}{quádrica} duplamente regrada.
  \item Classifique $x^2 + y^2 - 2x + 4y + 3 = 0$ em $\R^3$
        (complete os quadrados; identifique um cilindro
        circular reto e dê seu eixo e seu raio).
  \item Suponha agora $\operatorname{rank}A = 1$, digamos $\lambda_1
        \neq 0 = \lambda_2 = \lambda_3$. Rodando dentro do
        plano do núcleo e transladando, reduza a
        \[
        \lambda_1z_1^2 + 2\beta z_2 = 0 \quad (\beta \neq 0)
        \qquad\text{ou}\qquad
        \lambda_1z_1^2 + c'' = 0 :
        \]
        um cilindro parabólico, ou um par de planos paralelos, um
        plano duplo, ou o conjunto vazio. Classifique $(x + y)^2 =
        z$ completamente (forma normal, eixo de invariância por
        translação).
\end{enumerate}

\medskip
\noindent\textbf{Parte IV --- O teorema de classificação.}
\begin{enumerate}[resume]
  \item Reúna as Partes I--III num teorema: \emph{toda
        \omterm{pb:b2:surfaces:1}{quádrica} de $\R^3$ é levada por um movimento rígido sobre
        exatamente uma forma normal}. Liste os dezessete tipos
        \omterm{def:b2:affine:subspace}{afins} (conte as variantes vazias e os conjuntos degenerados)
        e destaque as nove \emph{superfícies} \omterm{pb:b2:surfaces:1}{quádricas}:
        elipsoide, hiperboloides de uma e de duas folhas, cone,
        paraboloides elíptico e hiperbólico, cilindros elíptico,
        hiperbólico e parabólico.
  \item Escreva o \emph{algoritmo} de classificação: dados $(A,
        b, c)$, que quantidades você calcula, em que
        \omterm{def:b2:structures:generated}{ordem}, e qual ramificação decide qual tipo? Justifique
        que cada passo é efetivo (\omterm{def:b2:reduction:eigen}{autovalores} de uma
        matriz $3\times3$ \omterm{def:b2:quadratic:adjoint}{simétrica}, posto, resolubilidade de
        $A\Omega = -b$).
  \item Rode o algoritmo em $x^2 + y^2 + z^2 - 2xy - 2yz -
        2zx = 1$: mostre que $A = 2I - J$ tem \omterm{def:b2:reduction:eigen}{espectro} $\{2, 2,
        -1\}$ e conclua: um hiperboloide de uma folha de
        revolução em torno de $\R(1,1,1)$.
  \item Rode-o em $x^2 + y^2 - z^2 - 2x + 4y + 2z + 4 = 0$:
        determine o centro e identifique a \omterm{pb:b2:surfaces:1}{quádrica}.
  \item Rode-o em $x^2 + 4xy + y^2 = 2z$: diagonalize o
        bloco $xy$ ($u = \frac{x+y}{\sqrt2}$, $v =
        \frac{x-y}{\sqrt2}$) e identifique a \omterm{pb:b2:surfaces:1}{quádrica}.
  \item Euclidiano contra \omterm{def:b2:affine:subspace}{afim}. Mostre que duas \omterm{pb:b2:surfaces:1}{quádricas}
        centrais em forma normal são \emph{rigidamente} equivalentes
        se e somente se têm as mesmas listas de coeficientes (a menos de
        permutação e de um escalar positivo comum sobre a
        equação), enquanto \emph{afimmente} só sobrevivem os
        dados de assinatura: todo elipsoide é imagem \omterm{def:b2:affine:subspace}{afim} da
        esfera redonda. Que teorema garante que a
        assinatura não pode mudar pelo caminho
        (\cref{thm:b2:quadratic:sylvester})?
\end{enumerate}

\medskip
\noindent\textbf{Parte V --- Dividendos.}
\begin{enumerate}[resume]
  \item Mostre que toda seção de uma \omterm{pb:b2:surfaces:1}{quádrica} por um plano
        \omterm{def:b2:affine:subspace}{afim} é uma cônica (possivelmente degenerada) desse plano.
        Identifique as seções da sela $z = xy$ pelos
        planos $z = c$ ($c \neq 0$ e $c = 0$).
  \item Que superfícies \omterm{pb:b2:surfaces:1}{quádricas} contêm retas? Mostre
        que o elipsoide, o hiperboloide de duas folhas e o
        paraboloide elíptico não contêm nenhuma \emph{(restrinja $q$
        a uma reta e use a desigualdade de Cauchy--Schwarz para
        o caso de duas folhas)}; que cone e cilindros são
        regrados por uma família; e conclua que as superfícies
        \omterm{pb:b2:surfaces:1}{quádricas} duplamente regradas são exatamente o hiperboloide
        de uma folha e o paraboloide hiperbólico.
  \item Quando se quer apenas o tipo \emph{\omterm{def:b2:affine:subspace}{afim}}, a
        \omterm{thm:b2:quadratic:gauss}{redução de Gauss} (\cref{thm:b2:quadratic:gauss}) é mais barata
        que diagonalizar. Refaça a questão 17 com o algoritmo
        de Gauss e confira a assinatura $(2, 1)$; que
        informação euclidiana Gauss perde?
  \item Todos os \omterm{def:b2:reduction:eigen}{autovalores} de uma matriz \omterm{def:b2:quadratic:adjoint}{simétrica} são reais; mostre
        que, em consequência, os \emph{sinais} dos \omterm{def:b2:reduction:eigen}{autovalores}
        de $A$ podem ser lidos no \omterm{def:b2:reduction:charpoly}{polinômio característico}
        pela regra dos sinais de Descartes, e verifique-o na
        questão 6: $\chi_A(\lambda) = \lambda^3 - 3\lambda^2
        - 9\lambda - 5$ tem exatamente uma mudança de sinal, logo
        assinatura $(1, 2)$.
  \item Síntese. Resuma o algoritmo em algumas linhas;
        enuncie o papel exato desempenhado (i) pelo teorema
        espectral, (ii) pela equação do centro $A\Omega = -b$,
        (iii) pelo teorema da inércia de Sylvester, (iv) pela \omterm{thm:b2:quadratic:gauss}{redução
        de Gauss}. O que a mesma máquina dá em
        $\R^2$, e o que muda em $\R^n$?
\end{enumerate}
\end{problem}
